\subsection{一次元固有側の双対切断と三次元AS正則性への条件付き還元}
\label{subsec:h3-order2-dual-hyperplane}
次は$(1,3)$と$(3,1)$に共通し、単冪部分を零と仮定しない。
全$(3,1)$の許容性を述べるものではなく、残る被覆を調べるための線型代数的還元である。

\begin{proposition}\label{prop:h3-order2-dual-hyperplane}
命題\ref{prop:h3-order2-mixed-rank}の一次元・三次元分解を
$B_1=U\oplus kz$、共通平面を$W\subset U$とする。
$E=B^!$、$\zeta=z^*\in E_1$、
$k\xi=W^\perp\subset U^*\subset E_1$と置く。
このとき
\[
 \zeta E=E\zeta,\qquad \zeta\xi=\xi\zeta=0,
 \qquad H_{E/(\zeta)}=(1+t)^3,
 \qquad H_{\operatorname{Ann}(\zeta)}=t+3t^2+3t^3+t^4.
\]
左・右annihilatorは等しい。
$F=E/(\zeta)$がFrobeniusであるための必要十分条件は
\[
 \operatorname{Ann}(\zeta)=E\xi=\xi E
\]
であり、この場合$\operatorname{Ann}(\xi)=(\zeta)$も成り立つ。
そのとき$F$はKoszulで、二次双対$F^!$は三生成元quadratic AS3であり、
自然写像$F^!\to B$は単射である。
Frobenius条件は、$F_1\times F_2\to F_3$と
$F_2\times F_1\to F_3$の二つの$3\times3$対合行列の行列式が非零となる条件である。
これを全$(3,1)$領域で自動的に満たすとはしない。
\end{proposition}
\begin{proof}
混合関係は$zw=f(w)z$、$w\in W$の二本である。
従って双対積$\zeta U^*$と$U^*\zeta$は同じ二次元の奇二次空間を張り、
両方の核は$k\xi$となる。さらに$\zeta^2\ne0$である。
実際、もし零なら四次topの$zz$行も零となり、
系\ref{cor:h3-order2-plane-quadric}の非退化$q_0$に反する。
従って$\zeta E_1=E_1\zeta$は三次元であり、生成元から
$\zeta E=E\zeta$を得る。
$E$のFrobenius対合で左乗法と右乗法を双対化すると、
$L_\zeta:E_i\to E_{i+1}$の階数は$i=0,1,2,3$でそれぞれ$1,3,3,1$である。
よって商とannihilatorのHilbert値は表示した通り。

$I=(\zeta)$は両側イデアルで、中山作用も$\zeta$をスカラー倍するから$I$を保存する。
Frobenius traceを$\lambda$とすると、
$\{a:\lambda(aI)=0\}$と$\{a:\lambda(Ia)=0\}$は中山公式で一致し、
非退化性と$I$の両側性により、それぞれ$I$の左・右annihilatorに等しい。
従って左右annihilatorは同じ両側イデアル$J$である。
$J$は$F$の双対に中山作用による捩れと次数移動を加えた両側加群である。
$F$がFrobeniusであることは、その双対が左右各々一元生成となることに同値。
$J_1=k\xi$なので、これは$J=E\xi=\xi E$に等しい。
この場合$\dim J=8$で、$\dim(\zeta)=8$でもある。
$\zeta\xi=\xi\zeta=0$と次元比較により
$\operatorname{Ann}(\xi)=(\zeta)$も得る。

従って$F$は$E$加群として、$\zeta$と$\xi$の乗法が交互に現れる線型自由解消
\[
 \cdots\longrightarrow E(-3)\longrightarrow E(-2)
 \longrightarrow E(-1)\longrightarrow E\longrightarrow F\longrightarrow0
\]
を持つ。$\operatorname{Tor}^E_q(F,k)$は内部次数$q$に集中し、
正次数の$F$作用は零である。
change-of-rings spectral sequence
\[
 \operatorname{Tor}^F_p\bigl(k,\operatorname{Tor}^E_q(F,k)\bigr)
 \ \Longrightarrow\ \operatorname{Tor}^E_{p+q}(k,k)
\]
を用いる。もし$F$がKoszulでなければ、最小の$p$で内部次数$j>p$の
$\operatorname{Tor}^F_p(k,k)$が存在する。
この成分を$q=0$行に置くと、入る高次微分は第一象限性により存在せず、
出る微分の行先は第一添字が$p$未満なので内部次数$p-1$に集中し、内部次数が合わない。
従ってこの成分は生き残り、Koszulな$E$のTorの対角集中性に反する。
よって$F$もKoszulである。Hilbert値$(1,3,3,1)$とFrobenius性から
二次双対$F^!$はAS3となる。
さらに今度はspectral sequenceの全$E_2$項が内部次数$p+q$に集中するので、
全ての高次微分は内部次数の不一致により零である。
従ってedge map
$\operatorname{Tor}^E_p(k,k)\to\operatorname{Tor}^F_p(k,k)$は全$p$で全射となる。
各次数を双対化すると
$\operatorname{Ext}^*_F(k,k)\to\operatorname{Ext}^*_E(k,k)$は単射である。
両辺のKoszul性による同定で、これはまさに自然な二次代数写像$F^!\to B$である。
最後の二行列判定は、$F_0,F_3$が一次元であることとFrobeniusの定義から従う。
\end{proof}

この命題による三生成元AS3部分代数の存在には、明記したFrobenius条件が必要である。
その非消滅を全$(3,1)$の被覆で無条件に仮定しない。
明示したAS4族上では$E,F$の各次数は局所自由で、上の行列式はFrobenius開部分を定める。
命題\ref{prop:h3-positive3-binormal-exclusion}の$a\ne0$族はこの開部分に属する。
実際、その双対では$\zeta=Z_*,\xi=P_*$、
$P_*^2=-aZ_*^2$、$Z_*^3=0$となり、
$F$は三生成元の量子外積代数、$\operatorname{Ann}(Z_*)=(P_*)$である。
\texttt{verify\_cubic\_order2\_dual\_hyperplane.py}はこの全係数模型の乗法、
Frobenius対合、四つの階数、相互annihilatorと商対合を厳密検算する。
一般のFrobenius非消滅と、その退化域でのH3許容性は課題として残す。
